% paper_macros.tex  ---  numeric macros used by paper.tex
%
% PROVENANCE.  This file is NOT fully regenerated by build.sh.  build.sh runs
% paper.py, whose _write_macros() emits only the core whole-sample macros; the
% values below are the manuscript-facing subset, each RECORDED from a named,
% tracked computed output and independently verified against it (see the audit
% trail in outputs/referee/ and referee_scripts/).  To update a value, re-run
% the cited analysis and copy the number here; do not edit by guesswork.
%
% Only macros actually referenced by paper.tex are kept (stale/unused macros
% were removed).  Every value here has been checked to match its source file.

% ---- Sample sizes (matched catalogs) ---------------------------------------
% TNG: SubLink-matched; SIMBA: spatially matched. Mid/low/high partition the
% TNG SubLink sample at logM* = 9.55 and 10.55 (2,902 + 3,430 + 1,030 = 7,362).
\newcommand{\TNGSubLinkN}{7{,}362}
\newcommand{\SIMBASpatialN}{4{,}203}
\newcommand{\MidMassN}{3{,}430}
\newcommand{\LowMassN}{2{,}902}
\newcommand{\HighMassN}{1{,}030}

% ---- Sliding-window scan  (outputs/baseline_B/results_mass_scan_l3.json) ----
% WindowNWindowsL = # windows with L3 internal marginal CI>0 (9.55-10.55);
% WindowPeakMarg = max L3 internal marginal (at centre 10.35).
\newcommand{\WindowLoL}{9.55}
\newcommand{\WindowHiL}{10.55}
\newcommand{\WindowNWindowsL}{11}
\newcommand{\WindowPeakMarg}{0.161}

% ---- Paired gap internal-halo under L3
%      (outputs/baseline_B/results_mass_scan_l3_quench_pg.json, paired_L3) -----
% NWindows = # windows with delta CI>0; Lo/Hi = min/max delta over 9.55-10.55.
\newcommand{\PairedGapNWindows}{13}
\newcommand{\PairedGapLo}{+0.052}
\newcommand{\PairedGapHi}{+0.089}

% ---- f_star / L3 absorption  (outputs/referee/fstar_absorption_summary.csv) -
\newcommand{\FstarInternalMarg}{0.086}
\newcommand{\FstarInternalCILo}{+0.043}
\newcommand{\FstarInternalCIHi}{+0.133}
\newcommand{\FstarHaloMarg}{0.017}
\newcommand{\LThreeInternalMarg}{0.135}
\newcommand{\LThreeHaloMarg}{0.083}

% ---- SIMBA L1 quenching peak (outputs/referee/simba_same_control_scores.csv)-
\newcommand{\SimbaQuenchPeak}{0.141}

% ---- Mid-mass raw family scores
%      (outputs/baseline_B/results_geoctrl_l2.json, delta_logmstar) -----------
\newcommand{\MidMassInternalR}{0.307}
\newcommand{\MidMassHaloR}{0.235}
\newcommand{\MidMassGap}{+0.072}

% ---- Spatial-vs-SubLink matching confound
%      (L1 internal baseline: baseline_A spatial vs baseline_B SubLink) -------
% SpatialInflation = SpatialBaselineTNG / SubLinkBaseline.
\newcommand{\SpatialInflation}{3.7}
\newcommand{\SubLinkBaseline}{0.059}
\newcommand{\SpatialBaselineTNG}{0.220}
\newcommand{\SpatialBaselineSIMBA}{0.273}

% ---- Assembly ablation
%      (outputs/referee/unresolved_absorption_feature_groups.csv;
%       outputs/baseline_B/results_geoctrl_l3_9p5_10p5_geomabl_*.json) --------
% Recovery = (ablated internal marginal) - (full-L3 internal marginal).
% Timing recovery is the source value 0.0253 (-> +0.025), not 0.161-0.135.
% Pct = recovery as a fraction of the published L3 decomposition.
\newcommand{\AblationTimingRecovery}{+0.025}
\newcommand{\AblationTimingPct}{19}
\newcommand{\AblationLongAccretionRecovery}{+0.014}
\newcommand{\AblationLongAccretionPct}{10}
\newcommand{\AblationMergerRecovery}{+0.000}
